\chapter{Segundo ensaio: custo-benefício e escolha do instrumento}
\label{sec:cb}

O primeiro ensaio mede o que a lei fez. O segundo faz duas perguntas sobre
esses números: se o benefício em saúde supera o custo em produção, e se outro
instrumento teria feito melhor.

\section{Custo-benefício}

Os benefícios são o efeito sobre a saúde ao nascer, valorado pela literatura de
retornos ao peso ao nascer e pelo custo hospitalar neonatal, e a variação nas
intoxicações, valorada pelo método do custo da doença de
\citeonline{soares2009socialcost, soares2012impactos}. Os valores publicados por
esses autores são do Paraná em 1998 e 1999 e servem de método; os preços serão
reestimados com as internações do Ceará e do Rio Grande do Norte. Os custos são
a variação de área, produção e emprego estimada no primeiro ensaio e a
diferença de custo entre as tecnologias de aplicação.

Dois cuidados. O benefício calculado é um piso, porque não inclui efeitos
crônicos nem o dano ambiental. E a substituição entra dos dois lados da conta:
se a troca do avião pela aplicação manual elevar a intoxicação de
trabalhadores, esse aumento é custo da política.

Sobre o custo há, hoje, duas afirmações em conflito e nenhuma estimativa com
contrafactual. Em 2010, na discussão da lei municipal, uma empresa previu perda
de metade da produção com a proibição \cite[p.~169]{marinho2011dimensao}. Uma
nota técnica mostra que a produção e a exportação de banana do Ceará subiram
depois de 2018 \cite{cavalcante2022nota}. A série de produção do IBGE confirma
a alta, e mostra também que ela parte do fim de uma seca de seis anos, cujo
piso foi 2016, e que a produção potiguar cresceu cerca de 40\% entre 2016 e
2018 e ficou estável depois. As trajetórias dos dois estados não eram paralelas
às vésperas da lei, o que é ameaça para o desfecho econômico e reforça o uso
dos municípios da Chapada, e não dos estados, como unidade.

O resultado será apresentado em cenários, com a pergunta invertida: que efeito
sobre o peso ao nascer seria necessário para a lei se pagar, dado o custo
estimado, e esse efeito está dentro ou fora do intervalo de confiança?

\section{O instrumento}

A teoria organiza a comparação. Sob incerteza, a regulação por quantidade supera
a por preço quando o dano marginal cresce mais depressa que o custo marginal de
evitá-lo \cite{weitzman1974}. A deriva é poluição difusa, sem medidor; uma
taxa sobre ela não é implementável, e uma taxa sobre o voo não distingue o
bananal ao lado da escola do bananal isolado. A regulação de agrotóxicos tem,
além disso, efeitos distributivos que dependem de quem consegue substituir
\cite{lichtenberg1988}.

Quatro fatos do caso brasileiro informam a discussão. A zona de exclusão já
existia na norma federal e era descumprida. O instrumento de preço opera com o
sinal trocado: a desoneração fiscal dos agrotóxicos somou R\$~9,86 bilhões em
2017 \cite{cunha2020incentivos}. A produtividade marginal estimada do
agrotóxico é baixa em quase todo o país e mais alta no Nordeste
\cite{rodrigues2022tese}; a estimativa é uma média regional, dominada por
agricultura de sequeiro, e sensível à forma funcional, mas sugere que o custo
de restringir o uso varia entre regiões. E
o legislador cearense já se moveu, em 2024, da proibição para a regulação do
método, ao admitir drones.

O ensaio registra de antemão como lerá cada resultado do primeiro
(Tabela~\ref{tab:instrumento}). As leituras são condicionadas ao que o
intervalo de confiança exclui, e não à ausência de significância.

\begin{table}[ht]
\centering\small
\caption{Leitura do instrumento conforme o resultado do primeiro ensaio}
\label{tab:instrumento}
\begin{tabularx}{\textwidth}{@{}LL@{}}
\toprule
Se o primeiro ensaio encontrar & Leitura \\
\midrule
O intervalo exclui o zero a favor de melhora ao nascer, e exclui perda relevante de produção & A proibição teve benefício com custo baixo; a discussão passa a ser se outro instrumento teria obtido o mesmo \\
O intervalo exclui o zero a favor de melhora ao nascer, e não exclui perda de produção & O custo-benefício decide, pelo efeito de equilíbrio calculado na seção anterior \\
O intervalo para nascimentos inclui o zero; o de intoxicações exclui o zero para cima & Há indício de troca de exposição ambiental por ocupacional; o intervalo para nascimentos diz qual ganho máximo é compatível com os dados \\
O primeiro estágio mostra que o avião já não voava em 2015--2018 & A lei consolidou um estado que a seca havia produzido; seu efeito é o de impedir a retomada, que o contraste de 2019 não mede \\
\bottomrule
\end{tabularx}
\fonte{Elaboração própria.}
\end{table}
